% TOP_PROBLEM: {"problemId": "problem.top500-omission-w045049-fpt-versus-w1", "canonicalTitle": "FPT versus W[1]", "releaseRank": 76, "primaryDomain": "theoretical_computer_science", "primaryDomainLabel": "Theoretical computer science", "displayStatus": "open", "releaseStatus": "open"}
% !TeX program = lualatex
% Definition and sources only; this file is not an LLM solution attempt.
\documentclass[11pt]{article}
\usepackage[margin=25mm]{geometry}
\usepackage{fontspec}
\setmainfont{Latin Modern Roman}
\usepackage{amsmath,amssymb,mathtools}
\usepackage{unicode-math}
\setmathfont{Latin Modern Math}
\usepackage{xurl,hyperref}
\hypersetup{colorlinks=true,urlcolor=blue}
\setcounter{secnumdepth}{0}
\setlength{\parindent}{0pt}
\setlength{\parskip}{5pt}
\setlength{\emergencystretch}{3em}
\begin{document}
\section*{\texorpdfstring{FPT versus W[1]}{FPT versus W[1]}}\label{76}
\subsection{Definitions and mathematical statement}
A parameterized decision problem is fixed-parameter tractable if it has a uniform algorithm taking time $f(k)n^c$, where $f$ is a total computable function and $c$ is independent of $k$. For the clique problem, $n=|V(G)|$ and a $k$-clique is a set of $k$ pairwise adjacent vertices. The equivalent concrete target is the nonexistence of
\[
 A,f,c\quad\text{with}\quad T_A(G,k)\le f(k)n^c
\]
for all finite simple $G$ and $1\le k\le n$. An $n^{g(k)}$ algorithm does not meet this condition. The conjectured separation is $\mathsf{FPT}\ne\mathsf{W[1]}$ in the standard uniform parameterized model.
\par\smallskip\textit{Formulation and concept references:} \href{https://arxiv.org/html/2304.07516v2}{[S1]}.

\subsection{Short English statement}
Can the combinatorial difficulty of finding a clique be confined entirely to its requested size, leaving a fixed polynomial dependence on graph size?
\subsection{Sources}
\begin{itemize}
\item[S1]Chen, Feng, Laekhanukit and Liu, Simple Combinatorial Construction, arXiv:2304.07516v2 (2024).
\url{https://arxiv.org/html/2304.07516v2}.
\textit{Formulation source.}
\item[S2]Sharp-W[1] = FPT: Fixed-Parameter Tractable Exact Algorithms for the \#k-Matching Problem.
\url{https://arxiv.org/abs/2604.16308}.
\textit{Further reference; not a proof endorsement.}
\item[S3]A Parameterized-Complexity Framework for Finding Local Optima.
\url{https://drops.dagstuhl.de/storage/00lipics/lipics-vol362-itcs2026/LIPIcs.ITCS.2026.66/LIPIcs.ITCS.2026.66.pdf}.
\textit{Further reference; not a proof endorsement.}
\end{itemize}
\end{document}
